\section{Architecture logicielle}
\label{sec:architecture}
\subsection{Une application modulaire à organisation MVC}
OctoGenere est une application Java unique, organisée en packages. Le terme \emph{module} désigne ici une séparation logique ; il ne s'agit ni de microservices ni de modules Maven déployés séparément. L'interface suit une organisation MVC adaptée à JavaFX : les vues collectent les actions et présentent les données ; \code{UIController} coordonne les traitements ; les services métier réalisent l'exploration, l'analyse et l'export.

\code{Main} est le point d'entrée unique déclaré dans Maven et délègue le démarrage à l'application JavaFX \code{MainApp}. Celle-ci constitue le point d'assemblage : elle charge la configuration et le catalogue, crée le fournisseur et le moteur, puis relie le contrôleur aux panneaux. Cette construction explicite rend les dépendances visibles sans imposer un framework d'injection. Le moteur dépend de \code{LlmProvider} et de \code{AnalysisObserver}, pas d'un contrôle graphique concret.

\begin{figure}[H]
\centering
\begin{tikzpicture}
\node[darkbox,text width=14.1cm] (views) at (0,0) {Vues JavaFX\\\footnotesize Sélection du projet · configuration · arborescence · résultats};
\node[box,text width=6.0cm] (controller) at (-3.8,-1.65) {\textbf{UIController}\\Actions, état et tâches asynchrones};
\node[box,text width=5.1cm] (explorer) at (4.15,-1.65) {\textbf{ProjectExplorerService}\\Lecture de l'arborescence};
\node[darkbox,text width=6.0cm] (engine) at (-3.8,-3.45) {\textbf{LlmEvaluationEngine}\\Orchestration de l'analyse};
\node[box,text width=5.1cm] (report) at (4.15,-3.45) {\textbf{GenerateReportUseCase}\\ReportGenerator / export LaTeX};
\node[box,text width=4.0cm,minimum height=1.5cm] (context) at (-5.1,-5.9) {\textbf{Contexte et consignes}\\XmlFileTreeBuilder\\PromptBuilder};
\node[box,text width=4.0cm,minimum height=1.5cm] (provider) at (0,-5.9) {\textbf{LlmProvider}\\GeminiProvider\\OpenAiProvider};
\node[box,text width=4.0cm,minimum height=1.5cm] (sandbox) at (5.1,-5.9) {\textbf{SandboxService}\\SandboxExecutor\\Docker};
\draw[flow] (views.south -| controller.north) -- (controller.north);
\draw[flow] (controller.east) -- (explorer.west);
\draw[flow] (controller.south) -- node[tag,right]{analyze} (engine.north);
\draw[flow] ([xshift=1.8cm]controller.south) -- ++(0,-0.35) -| (report.north);
\draw[flow] (engine.south) -- ++(0,-0.55) -| (context.north);
\draw[flow] (engine.south) -- ++(0,-0.55) -| (provider.north);
\draw[flow] (engine.south) -- ++(0,-0.55) -| (sandbox.north);
\end{tikzpicture}
\caption{Principales dépendances de service. Le contrôleur transmet à l'export le résultat métier renvoyé par le moteur ; le fournisseur est le seul composant reliant l'analyse au service LLM distant.}
\label{fig:architecture}
\end{figure}

\subsection{Données et responsabilités}
\code{EvaluationReport} regroupe l'identité du projet, la date, les paramètres descriptifs, le modèle, les résultats par critère et la synthèse. Chaque \code{CriterionResult} porte une note, son maximum et les observations associées. Ces records effectuent des contrôles de construction et des copies de collections, ce qui évite qu'une modification ultérieure d'une liste d'entrée altère un résultat déjà validé.

L'exploration visuelle relève de \code{ProjectExplorerService}, qui construit les nœuds de l'arbre. Le composant \code{XmlFileTreeBuilder} lit les contenus retenus pour le prompt. Ces deux parcours ont des finalités et des limites différentes ; leurs représentations ne sont donc pas strictement identiques.

\sourcecode{\code{Main.java}, \code{ui/MainApp.java}, \code{ui/controller/UIController.java}, packages \code{analysis/model}, \code{project} et \code{prompt/xml}}

\clearpage
\suite{Architecture logicielle — comportement dynamique}
\subsection{Une orchestration explicite de l'analyse}
Le moteur prépare le contexte, tente l'exécution Maven dans la sandbox, construit le prompt puis appelle le fournisseur. Les flux standard et d'erreur de la sandbox deviennent des observations jointes à la requête. Le LLM reçoit du texte et produit une réponse : il ne choisit pas de commande à lancer et ne dialogue pas directement avec Docker.

\begin{figure}[H]
\centering
\begin{tikzpicture}[x=1cm,y=1cm]
\node[darkbox,text width=2.2cm] (ui) at (0,0) {UIController};
\node[darkbox,text width=2.3cm] (eng) at (4.05,0) {Moteur d'analyse};
\node[box,text width=2.2cm] (sb) at (8.25,0) {Sandbox};
\node[box,text width=2.2cm] (llm) at (12.4,0) {Fournisseur LLM};
\foreach \x in {0,4.05,8.25,12.4} {\draw[dashed,line] (\x,-0.5) -- (\x,-8.05);}
\draw[flow] (0,-1) -- node[tag,above]{projet, critères, modèle} (4.05,-1);
\draw[flow] (4.05,-1.8) -- ++(0.65,0) -- ++(0,-0.45) -- (4.05,-2.25);
\node[tag,anchor=east] at (3.9,-2.0) {construction du contexte XML};
\draw[flow] (4.05,-2.85) -- node[tag,above]{commande Maven isolée} (8.25,-2.85);
\draw[backflow] (8.25,-3.55) -- node[tag,above]{sorties et code de retour} (4.05,-3.55);
\draw[flow] (4.05,-4.3) -- node[tag,above]{prompt : consignes, contexte et observations} (12.4,-4.3);
\draw[backflow] (12.4,-5.05) -- node[tag,above]{réponse textuelle au format JSON} (4.05,-5.05);
\draw[flow] (4.05,-5.65) -- ++(0.65,0) -- ++(0,-0.45) -- (4.05,-6.1);
\node[tag,anchor=east] at (3.9,-5.9) {validation des résultats};
\draw[backflow] (4.05,-6.85) -- node[tag,above]{EvaluationReport} (0,-6.85);
\node[tag,anchor=west,text width=10.7cm] at (0,-7.65) {Les notifications de progression sont émises pendant le traitement. L'export est une action ultérieure, déclenchée par l'utilisateur.};
\end{tikzpicture}
\caption{Scénario nominal d'une analyse. En cas de \code{SandboxException}, le moteur poursuit avec une mention d'indisponibilité ; une réponse LLM invalide interrompt la restitution des résultats.}
\label{fig:sequence}
\end{figure}

\subsection{Réactivité et cohérence de l'interface}
Le contrôleur exécute l'exploration, l'analyse et l'export dans un exécuteur dédié. Des objets \code{Task} organisent les retours de succès, d'échec et d'annulation. Les notifications métier sont replacées sur le thread JavaFX par \code{Platform.runLater}, conformément au modèle de concurrence de JavaFX \cite{javafx}. Un second exécuteur charge la liste des modèles indépendamment de ces opérations.

Pendant un traitement, les contrôles concernés sont désactivés. Le contrôleur ne conserve qu'un dernier rapport en mémoire et autorise son export après succès. Cette gestion empêche plusieurs actions concurrentes depuis les boutons, mais ne constitue pas un historique persistant. La synchronisation entre l'arrivée tardive de la liste des modèles et une analyse déjà lancée reste un cas à approfondir par des tests graphiques.

\clearpage
\section{Patrons de conception et justification}
\label{sec:patterns}
Les patrons sont retenus pour résoudre un problème concret de variation ou de couplage. Nous utilisons la terminologie classique des patrons objets \cite{gof}, en distinguant une implémentation complète d'une simple ressemblance structurelle.

\subsection{Adapter : isoler les protocoles externes}
\textbf{Problème.} Gemini et OpenAI n'exposent pas les mêmes classes de requête, structures de réponse et exceptions. Une dépendance directe du moteur envers ces API disperserait les particularités techniques dans la logique d'analyse.

\textbf{Réalisation.} \code{GeminiProvider} et \code{OpenAiProvider} réalisent \code{LlmProvider}. Ils construisent la requête propre au SDK, extraient le texte exploitable et traduisent les erreurs en \code{LlmException}. Le moteur peut donc appeler \code{ask} à travers un contrat commun. Le modèle choisi est un paramètre de la requête, distinct de l'identité du fournisseur.

\textbf{Intérêt et limite.} L'adaptation facilite le remplacement du SDK et les tests. Elle ne rend pas tous les modèles sémantiquement équivalents : formats supportés, quotas et capacités de génération restent dépendants du service.

\begin{figure}[H]
\centering
\begin{tikzpicture}
\node[darkbox,text width=4.3cm] (port) at (0,0) {\textit{« interface »}\\\textbf{LlmProvider}\\ask · availableModels};
\node[box,text width=4.3cm] (gem) at (-3.4,-2.2) {\textbf{GeminiProvider}\\SDK Google Gen AI};
\node[box,text width=4.3cm] (oai) at (3.4,-2.2) {\textbf{OpenAiProvider}\\SDK OpenAI};
\draw[dashed,-{Triangle[open,length=3mm,width=3mm]},accent] (gem.north) -- ++(0,0.45) -| ([xshift=-9mm]port.south);
\draw[dashed,-{Triangle[open,length=3mm,width=3mm]},accent] (oai.north) -- ++(0,0.45) -| ([xshift=9mm]port.south);
\end{tikzpicture}
\caption{Réalisation du contrat fournisseur par deux adaptateurs. Les flèches pointent vers l'interface implémentée.}
\end{figure}

\subsection{Fabrique simple : centraliser la création}
\code{LlmProviderFactory.from} lit \code{LLM\_PROVIDER} et construit l'adaptateur correspondant. La sélection est concentrée dans un \code{switch}, au lieu d'être répétée dans l'interface et le moteur. La fabrique répond à la question \emph{quel objet créer ?}, alors que l'adaptateur répond à \emph{comment utiliser ce service à travers notre contrat ?} Les deux mécanismes sont complémentaires.

Il s'agit d'une \textbf{fabrique simple}, pas d'une Abstract Factory construisant des familles d'objets, ni d'un Factory Method fondé sur la redéfinition d'une méthode par héritage. L'ajout d'un fournisseur nécessite encore de modifier ce point de création ; cette centralisation reste proportionnée à deux implémentations.

\subsection{Strategy : substituer un comportement}
\code{EvaluationEngine} permet au contrôleur de recevoir un moteur réel ou \code{DemoEvaluationEngine}. Ce dernier produit des résultats fictifs pour tester l'interface et l'export sans réseau. De même, \code{SandboxExecutor} permet à \code{SandboxService} de déléguer à Docker ou à un faux exécuteur. La substitution est réalisée par composition ; elle ne suppose pas un sélecteur de stratégie exposé à l'utilisateur.

\clearpage
\suite{Patrons de conception — notifications et construction}
\subsection{Observer : notifier sans dépendre de JavaFX}
\textbf{Problème.} Une analyse dure assez longtemps pour nécessiter un retour de progression, alors que le moteur doit rester utilisable hors de l'interface graphique.

\textbf{Réalisation.} \code{AnalysisObserver} expose \code{onProgressUpdate} et \code{onNewLogAdded}. Le contrôleur implémente ce contrat et fournit son instance à l'analyse. Le moteur émet les notifications sans connaître les panneaux. Le contrôleur assure ensuite leur restitution sur le thread de l'interface.

\textbf{Compromis.} Le mécanisme est une forme simple d'Observer, avec un observateur transmis par appel. Il ne comporte pas de registre d'abonnés ni de bus d'événements. Cette simplicité couvre le besoin actuel et facilite les tests : un observateur de test peut enregistrer les messages et vérifier que la progression atteint sa valeur finale.

\subsection{Facade : proposer une entrée de service limitée}
\code{SandboxService.run} offre une façade au sous-système d'exécution. L'appelant fournit un dossier et une commande ; il n'assemble pas les options de \code{docker run}, ne lit pas lui-même les flux du processus et ne gère pas son nettoyage. Ces détails sont pris en charge derrière \code{SandboxExecutor}.

La façade simplifie l'utilisation du sous-système ; \code{DockerSandboxExecutor} adapte ensuite l'outil en ligne de commande au contrat Java. Facade et Adapter interviennent donc à deux niveaux différents. Le choix des restrictions reste documenté et testable dans \code{DockerCommandBuilder} et \code{SandboxLimits}.

\subsection{Builder : construire une demande d'exécution cohérente}
\code{ExecutionRequest.Builder} associe deux éléments obligatoires, le répertoire et la commande, à des paramètres optionnels : limites, réseau et environnement. Il évite une multiplication de constructeurs difficiles à lire et produit une demande dont les listes et les associations sont copiées.

Ce Builder facilite l'évolution de la configuration sans obliger tous les appelants à renseigner chaque option. Il ne remplace pas une validation exhaustive de sécurité : les contraintes doivent toujours être appliquées par l'exécuteur. Les builders des SDK sont, quant à eux, des mécanismes fournis par ces bibliothèques, et non des patrons conçus par le groupe.

\subsection{Des abstractions volontairement limitées}
Les critères décrivent ici des libellés, des maxima et des paramètres d'activation. Une classe Strategy par critère aurait ajouté des objets sans encapsuler d'algorithme distinct. La solution retenue est un catalogue de données validé au démarrage, partagé avec le moteur et transmis à la vue.

L'arbre \code{FileNode} est une structure récursive : un dossier contient des enfants du même type. Il présente une parenté avec Composite, mais le code ne définit pas de hiérarchie explicite \emph{composant/feuille/composé} avec opérations polymorphes. De même, un enchaînement fixe d'étapes dans \code{LlmEvaluationEngine} ne suffit pas à constituer Template Method : aucune étape n'y est spécialisée par héritage.

\begin{pointcle}{Bilan de conception}
Le bénéfice des patterns est visible dans les tests et les points de substitution. La principale dette de couplage concerne la sandbox : le moteur crée encore directement sa façade par \code{createDefault()}, ce qui limite son remplacement au niveau du scénario d'analyse complet.
\end{pointcle}

\clearpage
